\documentclass[11pt]{article}
\usepackage[margin=1in]{geometry}
\usepackage{amsmath,amssymb,amsthm,mathtools}
\usepackage[T1]{fontenc}
\usepackage{lmodern}
\usepackage[colorlinks=true,linkcolor=blue,citecolor=blue,urlcolor=blue]{hyperref}
\newtheorem{theorem}{Theorem}
\newtheorem{lemma}[theorem]{Lemma}
\newtheorem{proposition}[theorem]{Proposition}
\theoremstyle{remark}\newtheorem{remark}[theorem]{Remark}
\newcommand{\TV}{d_{\mathrm{TV}}}
\newcommand{\Cov}{\operatorname{Cov}}
\newcommand{\Var}{\operatorname{Var}}
\newcommand{\Ent}{\operatorname{Ent}}
\newcommand{\diag}{\operatorname{Diag}}
\newcommand{\E}{\mathbb E}
\newcommand{\Prob}{\mathbb P}
\newcommand{\one}{\mathbf 1}
\newcommand{\zero}{\mathbf 0}
\newcommand{\sig}{\sigma}
\title{Near-linear Glauber mixing below the full\protect\\ferromagnetic field threshold}
\author{Anonymous}
\date{14 October 2026}
\begin{document}
\maketitle
\begin{center}\small Unrefereed candidate, version 2.0.0\\
\url{https://doi.org/10.5281/zenodo.23359097}\end{center}
\begin{abstract}
We prove worst-start mixing in
$O_{\beta,\gamma,\lambda}(n\log^2(2n)\log(1/\varepsilon))$
single-site heatbath updates for the heterogeneous ferromagnetic two-spin
systems of Feng--Guo--Yang throughout their regime
$\lambda<\lambda_c(\beta,\gamma)$. Degrees are unrestricted, and individual
edge and vertex parameters need not be bounded away from zero or infinity.
The argument combines their uniform influence theorem and computation-tree
representation with the scalar-field log--Sobolev criterion of Miao--Yin.
A small-field estimate makes the required covariance rate integrable.
Two coupon-collector windows then permit deterministic pinning of vertices
with large degree or bias. The remaining law has polynomially bounded
conditional marginals, while the original trajectory leaves the pinned
region with small probability. This removes the startup and marginal
dependence without asserting a uniform entropy inequality for the original
untrimmed law.
\end{abstract}

\section{Statement and relation to the source question}

All logarithms are natural. Fix positive constants $\beta,\gamma,\lambda$
such that
\begin{equation}\label{eq:triple}
 \beta\leq1<\gamma,\qquad \beta\gamma>1,\qquad
 \lambda<\lambda_c(\beta,\gamma):=
 (\gamma/\beta)^{\sqrt{\beta\gamma}/(\sqrt{\beta\gamma}-1)}.
\end{equation}
Let $G=(V,E)$ be any finite simple undirected graph, with $n=|V|\geq2$.
Assign finite positive parameters satisfying
\begin{equation}\label{eq:parameters}
 0<\lambda_v<\lambda,\qquad
 0<\beta_e\leq\beta,\quad \gamma_e\geq\gamma,\quad
 1<\beta_e\gamma_e\leq\beta\gamma.
\end{equation}
For $x\in\{0,1\}^V$, define
\begin{equation}\label{eq:measure}
 \mu(x)=\frac1Z
 \prod_{v:x_v=0}\lambda_v
 \prod_{\{u,v\}\in E:x_u=x_v=0}\beta_{uv}
 \prod_{\{u,v\}\in E:x_u=x_v=1}\gamma_{uv}.
\end{equation}
These are exactly the systems in \cite[Definition~4]{FGY26}.
One step of $P$ selects a uniformly random vertex independently of the
past and resamples its spin from its exact conditional law. Write
\[
 d(t)=\max_{z\in\{0,1\}^V}\TV(P^t(z,\cdot),\mu),\qquad
 t_{\mathrm{mix}}(\varepsilon)=\min\{t\geq0:d(t)\leq\varepsilon\}.
\]

\begin{theorem}\label{thm:main}
For each triple satisfying \eqref{eq:triple}, there is a finite constant
$K=K(\beta,\gamma,\lambda)$ such that every system
\eqref{eq:parameters}--\eqref{eq:measure} satisfies
\[
 t_{\mathrm{mix}}(\varepsilon)
 \leq\left\lceil K n\log^2(2n)\log(1/\varepsilon)\right\rceil,
 \qquad 0<\varepsilon\leq\tfrac12.
\]
The constant is uniform in the graph, its degrees, all admissible
heterogeneous parameters, and the initial state.
\end{theorem}

Feng--Guo--Yang prove near-linear mixing below
$\lambda_0=\sqrt{\gamma/\beta}$ and polynomial mixing throughout
$\lambda<\lambda_c$ \cite[Theorems~5--6]{FGY26}. Their introduction asks
whether near-linear worst-start mixing holds throughout this latter
regime. Theorem~\ref{thm:main} answers that question with logarithmic
exponent two. The bound counts single-site updates; it does not claim
constant arithmetic cost for evaluating a conditional probability on a
dense graph.

An earlier warm-start argument addresses startup \cite[Lemma~63]{FGY26}.
Here deterministic degree and bias trimming is combined with a
log--Sobolev estimate on the retained coordinates to improve the bound.
The non-elementary inputs are the influence and self-avoiding-walk
representation results of \cite{FGY26}, and the existing scalar-field
log--Sobolev criterion of \cite{MY26}. We state precisely the consequences
used below. The application does not prove a new general entropy theorem,
nor does it impose a marginal lower bound on the original law.

Two hypotheses prevent a direct substitution into the entropy criterion.
The original law has no uniform polynomial lower bound on occupied
marginals, and a constant covariance envelope alone gives a divergent
small-field stability integral. Lemma~\ref{lem:stability} repairs the
second issue; deterministic trimming and trajectory comparison repair
the first. The entropy-boosting alternative of Chen--Feng--Yin--Zhang
\cite[Definition~1.9 and Theorem~1.11]{CFYZ22} additionally requires
relative marginal stability and a pinned modified log--Sobolev bound.
Section~\ref{sec:obstructions} gives an explicit obstruction to a direct
degree-uniform application in the occupied-zero orientation.

\section{Parameterization and imported analytic tools}

Expanding the logarithm of \eqref{eq:measure}, up to an additive constant,
gives
\begin{equation}\label{eq:logistic}
 \log\mu(x)=\text{constant}+\sum_v\theta_vx_v+
 \sum_{\{u,v\}\in E}J_{uv}x_ux_v,
 \quad J_e=\log(\beta_e\gamma_e),\quad
 \theta_v=-\log\lambda_v-\sum_{e\ni v}\log\beta_e.
\end{equation}
In particular, with $J_*:=\log(\beta\gamma)>0$,
\begin{equation}\label{eq:logbounds}
 0<J_e\leq J_*,\qquad \theta_v\geq-\log\lambda.
\end{equation}
Thus the conditional probability of spin one is
$\sig(\theta_v+\sum_{u\sim v}J_{uv}x_u)$, where
$\sig(h)=(1+e^{-h})^{-1}$, and is at least
\begin{equation}\label{eq:q}
 q:=\frac1{1+\lambda}.
\end{equation}
The chain is attractive: common site selections and common update
uniforms preserve coordinatewise order. Notice that bounded $J_e$ follows
from the product bound in \eqref{eq:parameters}; individual $\gamma_e$
can still be arbitrarily large.

Throughout the analytic part, an \emph{occupied} coordinate means original
spin zero: $X_v=1-x_v$. Multiplying the weight by $s^{\sum_vX_v}$,
$0<s\leq1$, replaces every $\lambda_v$ by $s\lambda_v$.
Conditioning an original neighboring spin to zero multiplies the remaining
vertex fugacity by $\beta_e$; conditioning it to one multiplies that
fugacity by $1/\gamma_e$. Both multipliers are at most one. Consequently
arbitrary mixed pinnings and these scalar tilts preserve the edge
conditions and the field cap $s\lambda$.

\paragraph{Imported influence statements.}
For occupied indicators under any such law $\rho$, set
\[
 \eta(\rho):=\max_v\sum_{u\ne v}
 \left|\rho(X_v=1\mid X_u=1)-\rho(X_v=1\mid X_u=0)\right|.
\]
Feng--Guo--Yang's Theorem~24 supplies a constant $C_\infty$ depending
only on the fixed triple such that $\eta(\rho)\leq C_\infty$ for every
law in this class. We also use their computation-tree representation and
copy-influence comparison \cite[Proposition~14, Lemmas~27 and~29,
equation~(15)]{FGY26}:
incoming graph influence is bounded by twice the sum of the corresponding
tree influences over all levels. The tree recursion for zero-to-one odds
is
\begin{equation}\label{eq:tree}
 F_u(z_1,\ldots,z_d)=a_u\prod_{i=1}^d f_i(z_i),\qquad
 f_i(z)=\frac{\beta_i z+1}{z+\gamma_i}.
\end{equation}
Here $a_u$ includes absorbed fixed leaves and satisfies $0<a_u<s\lambda$.
The tree influences allow the selected boundary ratio to change from
zero to infinity; the other siblings of that selected boundary vertex
are restricted to $(0,\lambda)$ as in \cite[Definition~28]{FGY26}.
Other boundary assignments can vary with the selected vertex. These
details matter for the uniform estimate in the next section.

\paragraph{Imported log--Sobolev criterion.}
For a reversible kernel $P$ with stationary law $\rho$, our conventions are
\[
 \mathcal E_P(f,f)=\langle f,(I-P)f\rangle_\rho,
 \qquad \Ent_\rho(g)=\E_\rho[g\log g]-\E_\rho[g]\log\E_\rho[g],
\]
for $g\geq0$, with $0\log0=0$.
For uniform $m$-site heatbath dynamics,
$\mathcal E_P(f,f)=m^{-1}\sum_i\E_\rho[\Var_\rho(f\mid X_{-i})]$.
The ordinary log--Sobolev inequality bounds $\Ent_\rho(f^2)$ by this
quadratic form; it is not the modified inequality with
$\mathcal E_P(g,\log g)$. Time counts individual updates, not sweeps or
rate-one-per-site continuous time.
We use the following quantitative consequence of
\cite[Definition~1.1, Theorem~1.4 and Section~2.1]{MY26}.
Let $\rho$ be positive on a downward-closed family in $2^{[m]}$ with at
least two states. Suppose insertion ratios are at most $\kappa\geq1$,
occupied marginals under positive pinnings are at least $b>0$, and
the scalar tilted, positively pinned laws satisfy
\[
 \Cov(X)\preceq C(t)\diag(\E X),\qquad 0\leq t<1,
 \quad I=\int_0^1\frac{C(t)-1}{1-t}\,dt<\infty,
\]
where the scalar multiplier is $1-t$ and $C\geq1$ is nonincreasing.
Then the ordinary $m$-coordinate heatbath Dirichlet form satisfies
\begin{equation}\label{eq:MY}
 \Ent_\rho(f^2)\lesssim
 \kappa e^I m\log\frac{m\kappa}{b}\,\mathcal E(f,f).
\end{equation}
Its mixing bound is
\begin{equation}\label{eq:MYmix}
 t_{\mathrm{mix}}(\varepsilon)\lesssim
 \kappa e^I m\left[
 \log\frac{m\kappa}{b}\,
 \log\log\frac{e^2}{\rho_{\min}}+\log(1/\varepsilon)\right].
\end{equation}
Proposition~2.1 of \cite{MY26} also gives
$\operatorname{gap}^{-1}\leq m(1+\kappa)e^I$ without a marginal lower
bound. These are imported results, including their ordinary update
normalization. Only the entropy estimate needs trimming in our application.

\section{An integrable scalar covariance rate}

\begin{lemma}\label{lem:stability}
There is $A=A(\beta,\gamma,\lambda)<\infty$ such that every law obtained
from \eqref{eq:measure} by arbitrary mixed pinning and multiplication by
$s^{\sum X_v}$ satisfies
\[
 \Cov(X)\preceq(1+As)\diag(\E X),\qquad 0<s\leq1,
\]
on its unpinned coordinates. In particular the rate $C(t)=1+A(1-t)$
has stability integral $I=A$.
\end{lemma}
\begin{proof}
Put
\[
 q_*:=\frac{\lambda+1}{\lambda+\gamma}<1,\qquad
 R:=\frac{\lambda}{(1-q_*)^2},\qquad
 s_0:=\min\{1,(2R)^{-1}\}.
\]
For all nonnegative $z$, the factors in \eqref{eq:tree} satisfy
$f_i(z)\leq1$. Hence every unpinned tree ratio is at most $a_u<s\lambda$,
even when its descendants have boundary ratios zero or infinity.
For $0\leq z\leq\lambda$, monotonicity gives
\[
 f_i(z)\leq q_*,\qquad
 0\leq f_i'(z)=\frac{\beta_i\gamma_i-1}{(z+\gamma_i)^2}
 \leq\frac{\gamma_i-1}{\gamma_i^2}\leq1.
\]
It follows that, when all child ratios lie in $[0,\lambda]$,
\begin{equation}\label{eq:derivative}
 0\leq\partial_iF_u\leq s\lambda q_*^{d-1}.
\end{equation}
At distance one, the exact endpoint difference with child $i$ changing
from zero to infinity is
\[
 a_u(\beta_i-1/\gamma_i)\prod_{j\ne i}f_j(z_j)
 \leq s\lambda q_*^{d-1},
\]
because the other siblings are below $\lambda$ in the imported tree
influence definition. Summing over children costs at most
$s\lambda d q_*^{d-1}\leq Rs$.

For distance $j\geq2$, the immediate child ratios are unpinned ratios,
so both endpoint values and their interpolation lie in $[0,\lambda]$.
Apply \eqref{eq:derivative} along the affected branch. Its bound is
independent of all boundary assignments, including assignments chosen
separately for different affected vertices. If $S_j(u)$ is the sum of
level-$j$ influences at $u$, this gives
\[
 S_1(u)\leq Rs,\qquad
 S_j(u)\leq Rs\max_{z\text{ child of }u}S_{j-1}(z),
 \qquad S_j(u)\leq(Rs)^j.
\]
Empty levels contribute zero. The imported graph-to-tree comparison
therefore yields, for $Rs<1$,
\[
 \eta(\rho)\leq2\sum_{j\geq1}(Rs)^j
 =\frac{2Rs}{1-Rs}.
\]
For $s\leq s_0$ this is at most $4Rs$. For $s\geq s_0$ use the fixed-cap
bound $C_\infty$ instead. Thus
\begin{equation}\label{eq:eta}
 \eta(\rho)\leq As,\qquad
 A:=\max\{4R,C_\infty/s_0\}.
\end{equation}
The restriction $\lambda<\lambda_c$ enters through $C_\infty$;
no small-field potential constant has been assumed uniform.

For completeness, the normalization conversion is as follows.
Write $p_i=\E_\rho X_i$, $V=\diag(p_i(1-p_i))$ and
$\Sigma=\Cov_\rho(X)$. Full support gives $0<p_i<1$ for unpinned sites.
The matrix $B=V^{-1}\Sigma$ has diagonal one and off-diagonal entries
\[
 B_{ij}=\rho(X_j=1\mid X_i=1)-\rho(X_j=1\mid X_i=0).
\]
Every absolute column sum is at most $1+As$ by \eqref{eq:eta}.
The symmetric positive semidefinite matrix $V^{-1/2}\Sigma V^{-1/2}$
is similar to $B$, so its largest eigenvalue is bounded by the induced
column norm of $B$. Congruence gives
\[
 \Sigma\preceq(1+As)V\preceq(1+As)\diag(p).
\]
This is a spectral-radius argument, not an assertion that similarity
preserves Euclidean operator norms. If no coordinates remain the condition
is vacuous. Finally $\int_0^1 A(1-t)/(1-t)\,dt=A$.
\end{proof}

\section{Startup from an arbitrary configuration}

The next estimate uses fresh uniforms selected by the site schedule.
It does not require the neighbors of a vertex to be independent in the
Gibbs law or in the evolving chain.

\begin{lemma}\label{lem:update}
Condition only on a complete deterministic site-selection schedule.
At an update of $v$ for which every neighbor has updated earlier,
the probability that the new spin is zero is at most
\[
 \lambda\rho^{\deg_G(v)},\qquad
 \rho:=(1-q)\beta+q/\gamma<1.
\]
At any update of $v$, its zero probability is also at most $e^{-\theta_v}$.
Both statements hold for every initial configuration.
\end{lemma}
\begin{proof}
For each neighbor $u$, select its last update uniform before this update
of $v$. Conditional on the schedule, their time indices are fixed,
distinct, and earlier than the fresh uniform at $v$. Thus their
threshold-$q$ indicators $B_u$ are independent Bernoulli$(q)$ variables.
Equation~\eqref{eq:q} implies $x_u\geq B_u$ pointwise at the time in
question. The actual spins may depend on many of these uniforms; their
independence is neither asserted nor needed.

The conditional zero odds at $v$ are
\[
 \lambda_v\prod_{u\sim v}\beta_{uv}^{1-x_u}\gamma_{uv}^{-x_u}.
\]
Each factor decreases with $x_u$ since $\beta_e\gamma_e>1$, and zero
probability is at most zero odds. Domination by the $B_u$ and independence
of the selected uniforms give
\[
 \Prob(x_v\text{ becomes }0\mid\text{schedule})
 \leq\lambda\prod_{u\sim v}[(1-q)\beta_{uv}+q/\gamma_{uv}]
 \leq\lambda\rho^{\deg_G(v)}.
\]
Here $\rho<1$ because $\beta\leq1$ and $1/\gamma<1$.
The second bound follows directly from
$1-\sig(\theta_v+\sum J_{uv}x_u)\leq e^{-\theta_v}$.
We have conditioned on no uniform values or realized neighbor spins.
\end{proof}

Fix the numerical exponent $a=10$ and set
\begin{equation}\label{eq:DB}
 D=\frac{a+2}{-\log\rho},\qquad B=a+2.
\end{equation}
Let $H\geq\max\{2,\lambda\}$ be a parameter-only constant to be made
larger at the end, and define
\begin{equation}\label{eq:trim}
 L=\left\lceil(a+3)n\log(Hn)\right\rceil,\quad
 T=\{v:\deg_G(v)>D\log(Hn)\text{ or }\theta_v>B\log(Hn)\},
 \quad \delta=(Hn)^{-a}.
\end{equation}
The set $T$ is determined by the input parameters, independently of all
chain randomness.

\begin{lemma}\label{lem:burn}
For all $v\in T$, all $t\geq2L$, and all initial states $z$,
\[
 \Prob_z(x_v(t)=0)\leq\delta,
 \qquad \mu(x_v=0)\leq\delta.
\]
Consequently, with $\widehat\nu=\mu(\,\cdot\mid x_T=\one)$,
$\TV(\mu,\widehat\nu)\leq n\delta$.
\end{lemma}
\begin{proof}
Two successive windows of $L$ independent uniform site selections both
cover $V$ except with probability at most
\[
 2n e^{-L/n}\leq2n(Hn)^{-(a+3)}\leq(Hn)^{-(a+1)}.
\]
On this schedule event, the last update of each $v$ in the second window
occurs after every neighbor updated in the first. The conditioning in
Lemma~\ref{lem:update} remains valid, because the event and the last-update
indices depend only on the schedule. For a high-degree vertex,
\[
 \lambda\rho^{\deg_G(v)}\leq\lambda(Hn)^{-(a+2)}
 \leq(Hn)^{-(a+1)}.
\]
For a large-bias vertex, $e^{-\theta_v}\leq(Hn)^{-(a+2)}$.
Adding the schedule failure shows, from every start,
$\Prob_z(x_v(2L)=0)\leq2(Hn)^{-(a+1)}\leq\delta$.

Let $\xi_t$ denote the law started from $\zero$. The minimality of
$\zero$ implies $\xi_0\preceq\xi_1$ in stochastic order; attractiveness
then implies $\xi_t\preceq\xi_{t+1}$ for every $t$. Every other starting
law is stochastically above $\xi_t$ at time $t$. Since $\{x_v=0\}$ is a
decreasing event, the bound at $2L$ therefore holds at every later time
from every start. The finite full-support heatbath chain is irreducible
and aperiodic. Taking its limit from $\zero$ transfers the bound to $\mu$,
without assuming any convergence rate. Finally the distance from a law
to its conditioning on an event equals the probability of the complement;
a union bound gives the asserted total variation estimate.
\end{proof}

\section{The residual law and its mixing time}

Let $U=V\setminus T$, $m=|U|$, and let $\nu$ be the marginal on $U$ of
$\widehat\nu$. Its effective logistic biases satisfy
\[
 -\log\lambda\leq\theta_v^{\mathrm{eff}}
 =\theta_v+\sum_{u\in T:u\sim v}J_{uv}
 \leq(B+J_*D)\log(Hn),\qquad v\in U.
\]
The original degree of each retained vertex is at most $D\log(Hn)$.
Thus, under any further configuration or pinning of retained coordinates,
the conditional logistic field is at most $M\log(Hn)$, where
\begin{equation}\label{eq:b}
M:=B+J_*D,\qquad b:=\frac1{1+(Hn)^M}\geq\frac12(Hn)^{-M}.
\end{equation}
The conditional zero probability is at least $b$ and the conditional one
probability is at least $q$. To obtain $M$, count pinned and unpinned
neighbors together: their combined contribution is at most
$J_*\deg_G(v)$, so no neighbor is counted twice.
Sequentially exposing all coordinates proves
\begin{equation}\label{eq:minmass}
 \nu_{\min}\geq\min\{b,q\}^{m},\qquad
 \log(1/\nu_{\min})=O_{\beta,\gamma,\lambda}(m\log(Hn)).
\end{equation}

In occupied coordinates the support is the full cube, hence downward
closed. An insertion ratio, holding all other spins fixed, is the
conditional odds of original zero and is at most $\lambda$:
every factor $\beta_e$ or $1/\gamma_e$ is at most one.
We may therefore use $\kappa=\max\{1,\lambda\}$. The lower bound $b$
holds even under mixed pinnings and in particular under the positive
occupied pinnings required by \eqref{eq:MY}. Lemma~\ref{lem:stability}
applies after the original-one pins $T$, so the integral is $I=A$.
All assumptions of the imported criterion have now been verified for
the residual law.

For $m\geq2$, \eqref{eq:MY} gives inverse log--Sobolev constant
$O_{\beta,\gamma,\lambda}(m\log(Hn))$.
Equations~\eqref{eq:MYmix}--\eqref{eq:minmass} imply constant-accuracy
mixing in $O_{\beta,\gamma,\lambda}(m\log^2(Hn))$ active updates:
indeed $m\leq n$, and both
$\log(m\kappa/b)$ and $\log\log(e^2/\nu_{\min})$
are $O_{\beta,\gamma,\lambda}(\log(Hn))$.

We need the chain which still selects among all $n$ original vertices
and does nothing on selecting $T$. Its kernel on $U$ is
\[
 Q=\frac mn P_U+\left(1-\frac mn\right)I.
\]
Its Dirichlet form is $m/n$ times that of $P_U$, so both the
log--Sobolev and spectral-gap constants rescale by $m/n$.
Heatbath kernels are averages of conditional-expectation projections
and have nonnegative spectrum; the same is true of $Q$. Applying the
same log--Sobolev mixing estimate gives a constant
$C_0=C_0(\beta,\gamma,\lambda)\geq1$, independent of $H,n$, such that
\begin{equation}\label{eq:coretime}
 t_*:=\left\lceil C_0n\log^2(Hn)\right\rceil
 \quad\Longrightarrow\quad
 \max_y\TV(Q^{t_*}(y,\cdot),\nu)\leq\tfrac1{16}.
\end{equation}
All dependence on $H$ enters through the displayed logarithms, so $C_0$
can be fixed before increasing $H$. This rescaling does not replace a
random number of active updates by its expectation.

If $m=0$, the extended residual law is a point mass and its mixing error
is zero. If $m=1$, the retained spin is refreshed at each selection of
its vertex, and the distance is at most $(1-1/n)^t\leq e^{-t/n}$.
Increasing the same parameter-only $C_0$ if necessary gives
\eqref{eq:coretime} also in this case. Disconnected graphs and isolated
vertices require no separate assumptions.

\section{Comparison, choice of constants, and arbitrary accuracy}

The following elementary comparison isolates the reusable coupling step.
\begin{lemma}[Finite-horizon pinning comparison]\label{lem:transfer}
Let $\mu$ have full support on a finite product space, and let $P$ be
heatbath dynamics with fixed state-independent site-selection probabilities.
Fix a deterministic pinning event $F=\{x_T=a_T\}$ and set
$\nu=\mu(\cdot\mid F)$. Let $Q$ use the original clock, ignore selections
in $T$, and update the other coordinates according to $\nu$. Then for
integers $b,t\geq0$,
\[
 d_P(b+t)\leq d_Q(t)+
 \sup_z\Prob_z(\exists j\in\{0,\ldots,t\}:X_{b+j}\notin F)+\mu(F^c),
\]
where both distances are worst-start distances to the respective stationary
laws. No attractiveness assumption is required for this comparison.
\end{lemma}
\begin{proof}
At time $b$, retain the original chain's unpinned coordinates and replace
its pinned coordinates by $a_T$. Full support makes this an admissible
state. Use fresh common site choices and heatbath randomness thereafter.
The chains agree while the original chain remains in $F$, because their
conditional update laws outside $T$ then coincide. The terminal coupling
error is bounded by the displayed exit probability, including time zero.
The pinned initial law is a mixture of states in $F$, so its mixing error
is at most $d_Q(t)$. Finally $\TV(\nu,\mu)=\mu(F^c)$; apply the triangle
inequality. No independence of successive exit events is used.
\end{proof}

\begin{proof}[Proof of Theorem~\ref{thm:main}]
Run the original chain from an arbitrary deterministic state until time
$2L$. Initialize a pinned chain by retaining its $U$-coordinates and
setting all $T$-coordinates to one. Thereafter use common site selections
and uniforms; the pinned chain ignores selections in $T$. The future
randomness is fresh. Although its initial residual state may be correlated
with the burn-in, \eqref{eq:coretime} holds for every state and hence every
mixture of such states.

If all original pins are one initially and stay one throughout the next
$t_*$ steps, the coupled configurations agree throughout: updates in $U$
have identical conditional laws, and updates in $T$ preserve their
common value on this event. Lemma~\ref{lem:burn} and a union bound over
vertices and integer times bound the failure probability by
$(t_*+1)n\delta$. This is a union bound on the original trajectory,
not an estimate of a hazard conditioned on avoiding an exit.
Together with the stationary conditioning error it gives
\begin{equation}\label{eq:compare}
 d(2L+t_*)\leq\frac1{16}+(t_*+2)n(Hn)^{-a}.
\end{equation}

The order of choices is important. The fixed triple determines
$q,\rho,J_*,D,B,M,A,\kappa$ and then $C_0$, all independently of $H$.
For $H,n\geq2$, $\log^2(Hn)\leq Hn$ and $\log^2(Hn)\geq1$, so
\[
 (t_*+2)n(Hn)^{-a}
 \leq(C_0+3)n^2\log^2(Hn)(Hn)^{-a}
 \leq(C_0+3)H^{1-a}n^{3-a}
 \leq(C_0+3)H^{1-a}.
\]
Now choose $H\geq\max\{2,\lambda\}$ large enough that the last expression
is at most $1/16$. This is possible because $a=10$ and $C_0$ does not
depend on $H$. Equation~\eqref{eq:compare} yields
\[
 d(t_0)\leq\tfrac18,\qquad
 t_0:=2L+t_*=O_{\beta,\gamma,\lambda}(n\log^2(2n)).
\]
The last estimate uses that $H$ is now a fixed parameter-only constant.

Define the diameter
$\bar d(t)=\max_{x,y}\TV(P^t(x,\cdot),P^t(y,\cdot))$.
The elementary contraction coefficient inequality gives
$\bar d(s+t)\leq\bar d(s)\bar d(t)$, and stationarity and the triangle
inequality give $d(t)\leq\bar d(t)\leq2d(t)$.
Thus $\bar d(t_0)\leq1/4$ and $d(jt_0)\leq4^{-j}$ for every integer
$j\geq1$. Taking
\[
 j=\left\lceil\frac{\log(1/\varepsilon)}{\log4}\right\rceil
\]
proves the result after absorbing ceilings into $K$, since
$0<\varepsilon\leq1/2$. The pins and the constant $H$ were chosen once,
independently of $\varepsilon$.
\end{proof}

\section{Linear relaxation and obstructions to direct entropy bounds}
\label{sec:obstructions}

\begin{proposition}\label{prop:gap}
With $A$ from Lemma~\ref{lem:stability} and
$\kappa=\max\{1,\lambda\}$, the original untrimmed chain satisfies
\[
 \operatorname{gap}(P)^{-1}\leq n(1+\kappa)e^A.
\]
The same conclusion holds after any mixed pinning, with $n$ replaced
by the number of remaining coordinates (when nonzero).
\end{proposition}
\begin{proof}
In occupied coordinates the law has full-cube support. Its insertion
ratios are at most $\lambda$, since each neighboring factor is either
$\beta_e$ or $1/\gamma_e$, both at most one.
Lemma~\ref{lem:stability} supplies the scalar-field stability integral
$I=A$ under every required positive pinning. Apply
\cite[Proposition~2.1]{MY26}. Mixed pinnings preserve the class.
\end{proof}

This is a consequence of an imported gap comparison, not a replacement
for the worst-start argument: the usual gap estimate retains dependence
on the smallest stationary mass.

For fixed state-independent site probabilities $p_v>0$ with
$\sum_vp_v=1$, the corresponding kernel $P_p$ also satisfies
\[
 \operatorname{gap}(P_p)^{-1}\leq
 \frac{(1+\kappa)e^A}{\min_vp_v}.
\]
Indeed $\mathcal E_{P_p}(f,f)\geq n(\min_vp_v)\mathcal E_P(f,f)$,
so the Poincar\'e variational principle gives the claim. This comparison
alone does not establish a nonuniform-scan worst-start mixing bound.

For a concrete obstruction to a uniform ordinary log--Sobolev constant,
take two isolated vertices in the envelope $(\beta,\gamma,\lambda)=(1,2,3)$,
with $\mu(x_1=0)=p\downarrow0$. For
$f=\one_{\{x_1=0\}}$,
\[
 \Ent_\mu(f^2)=p\log(1/p),\qquad
 \mathcal E_P(f,f)=\tfrac12p(1-p).
\]
Thus the inverse ordinary log--Sobolev constant is at least
$2\log(1/p)/(1-p)$ and diverges. Nevertheless a common-refresh coupling
gives $d(t)\leq2^{1-t}$ for $t\geq1$, uniformly in $p$, because the
chains agree once both sites have been selected.

For the relative-marginal-stability obstruction, take a star with $d$
leaves, $\beta_e=1$, $\gamma_e=2$ and all $\lambda_v=1$, again within
the envelope $(1,2,3)$. Summing independently over the leaves conditional
on the center gives center zero-to-one odds $(2/3)^d$.
Pinning all leaves to original zero changes those odds to $1$.
The ratio is $(3/2)^d$, so the constant in
\cite[Definition~1.9]{CFYZ22} is not degree-uniform in this orientation.
This excludes a direct uniform application of that criterion, not other
entropy arguments. Both examples were suggested by the supplied review.

\section{Scope of the conclusion}

The proof handles arbitrarily small positive $\lambda_v$ and $\beta_e$,
arbitrarily large individual $\gamma_e$, and all finite degrees allowed
by \eqref{eq:parameters}. Finite positivity ensures full support and
ordinary finite-chain convergence; no uniform minimum mass is assumed
before trimming. The finite triple remains strictly below $\lambda_c$;
the constants may deteriorate as the triple approaches its boundary.

The scalar estimate uses Feng--Guo--Yang's established correlation
theorem, and the entropy conversion is Miao--Yin's established theorem.
The role of deterministic trimming is to supply the latter theorem's
marginal and startup bounds on a residual law while controlling the
original dynamics from every start. In particular, the argument asserts
neither a degree-free uniform log--Sobolev constant for the original law
nor mixing beyond the stated parameter class. A generic change of scan
order or a bipartite lift is not used.

The reusable step is the finite-horizon comparison: a deterministic set
whose adverse coordinates remain rare after burn-in can be pinned for
analysis, provided the resulting law mixes quickly and the accumulated
trajectory and stationary-conditioning errors are controlled. We do not
claim this elementary coupling principle as a new general theorem.
The sharper $O(n\log n)$ bound and the threshold endpoint remain open
here; additional routine corollaries would not resolve either question.

\paragraph{Availability and provenance.}
The accompanying package contains the source, explicit claim boundaries
and reproducible finite checks. Those checks support selected identities,
not the universal analytic theorem or the imported source theorems.
This anonymous manuscript was developed and revised with AI assistance;
no external peer review or formal verification is claimed. No human
research contribution, external funding or conflict is asserted without
supporting records. This theoretical work involves no human or animal
subjects.

\begin{thebibliography}{FGY26}
\bibitem[FGY26]{FGY26}
Weiming Feng, Heng Guo, and Minji Yang.
\newblock \emph{Rapid mixing in positively weighted restricted Boltzmann machines}.
\newblock arXiv:2604.00963v2, 2026.
\newblock \url{https://arxiv.org/abs/2604.00963v2}.

\bibitem[MY26]{MY26}
Tianshun Miao and Yitong Yin.
\newblock \emph{Sampling Matchings in Near-linear Time}.
\newblock arXiv:2609.21936v1, 2026.
\newblock \url{https://arxiv.org/abs/2609.21936v1}.
\bibitem[CFYZ22]{CFYZ22}
Xiaoyu Chen, Weiming Feng, Yitong Yin, and Xinyuan Zhang.
\newblock \emph{Optimal mixing for two-state anti-ferromagnetic spin systems}.
\newblock arXiv:2203.07771v1, 2022.
\newblock \url{https://arxiv.org/abs/2203.07771v1}.
\end{thebibliography}
\end{document}
